\section{Resultados}
\label{sec:resultados}

\subsection{Tablas}

Las tablas usan \texttt{booktabs} (sin líneas verticales). Las columnas numéricas con tipo
\texttt{S} de \texttt{siunitx} alinean por la coma decimal (\cref{tab:manual}).

\begin{table}[htbp]
  \centering
  \caption{Comparación de modelos (validación cruzada de 10 particiones).}
  \label{tab:manual}
  \begin{tabular}{l S[table-format=1.3] S[table-format=1.3] S[table-format=1.3]}
    \toprule
    \textbf{Modelo} & {\textbf{Precisión}} & {\textbf{Recall}} & {\textbf{F1}} \\
    \midrule
    Regresión logística & 0.801 & 0.765 & 0.783 \\
    Árbol de decisión   & 0.832 & 0.790 & 0.810 \\
    Random forest       & \bfseries 0.874 & \bfseries 0.842 & \bfseries 0.858 \\
    \bottomrule
  \end{tabular}
\end{table}

Si los resultados se generan con Python, se pueden exportar a CSV y leerlos directamente con
\verb|\tablacsv|, sin copiar números a mano (\cref{tab:csv}).

\begin{table}[htbp]
  \centering
  \caption{Métricas leídas desde \texttt{datos/resultados.csv}.}
  \label{tab:csv}
  \tablacsv{datos/resultados.csv}
\end{table}

\subsection{Gráficos}

Con \texttt{pgfplots} los gráficos se dibujan desde un CSV y comparten la tipografía del
documento (\cref{fig:codo}).

\begin{figure}[htbp]
  \centering
  \begin{tikzpicture}
    \begin{axis}[width=0.7\textwidth, height=55mm, xlabel={Número de grupos $k$},
                 ylabel={Inercia intra-grupo}, grid=major, grid style={black!10},
                 xtick={1,...,8}, ymin=0]
      \addplot[acento, thick, mark=*] table[x=k, y=inercia, col sep=comma]{datos/curva_codo.csv};
      \draw[BrickRed, dashed] (axis cs:4,0) -- (axis cs:4,1600);
    \end{axis}
  \end{tikzpicture}
  \caption{Método del codo para elegir $k$.}
  \label{fig:codo}
\end{figure}

\subsection{Matriz de confusión}

\begin{table}[htbp]
  \centering
  \caption{Matriz de confusión del modelo seleccionado.}
  \label{tab:confusion}
  \begin{tabular}{cc|cc}
    & & \multicolumn{2}{c}{\textbf{Predicción}} \\
    & & Positivo & Negativo \\
    \midrule
    \multirow{2}{*}{\rotatebox{90}{\textbf{Real}}}
      & Positivo & \cellcolor{acento!25} 421 & 79 \\
      & Negativo & 61 & \cellcolor{acento!25} 939 \\
  \end{tabular}
\end{table}

Los números con unidades se escriben con \verb|\num| y \verb|\SI|:
\num{1234567.891}, \SI{3.5}{\giga\byte}, \SI{250}{\milli\second}.
